\section{有限秩自同态的迹}

在本小节中，A 表示一个交换环。设 E 是一个投射 A-模。这时 \(\operatorname{End}_{\mathrm{A}}^{\mathrm{f}}(\mathrm{E})\) 是 \(\operatorname{End}_{\mathrm{A}}(\mathrm{E})\) 的一个 A-子模，且典范映射 \(E^{*}\otimes_{A}E\to\operatorname{End}_{\mathrm{A}}(\mathrm{E})\) 定义 A-模的一个同构 \(\theta_{E}\)（从 \(E^{*}\otimes_{A}E\) 到 \(\operatorname{End}_{\mathrm{A}}^{\mathrm{f}}(\mathrm{E})\)）（引理 1）。考虑典范线性型 \(\tau:E^{*}\otimes_{A}E\to A\)（II, §4, No.~3, 第 273 页），它由公式 \(\tau(x^{*}\otimes x)=\langle x,x^{*}\rangle\) 刻画。把它与同构 \(\theta_{E}^{-1}\) 复合，我们得到一个线性型 \(\operatorname{Tr}: \operatorname{End}_{\mathrm{A}}^{\mathrm{f}}(\mathrm{E})\to\mathrm{A}\)，称为迹型。当 A-模 E 有限生成时，我们重新得到 II, §4, No.~3, 第 273 页中的定义。

命题 1. --- 设 E 是一个自由 A-模。设 \((e_{i})_{i\in\mathrm{I}}\) 是 E 的一个基，\((e_{i}^{*})_{i\in\mathrm{I}}\) 是它的对偶基。设 \(u\in\operatorname{End}_{\mathrm{A}}^{\mathrm{f}}(\mathrm{E})\)。族 \((\langle u(e_{i}),e_{i}^{*}\rangle)_{i\in\mathrm{I}}\) 有有限支撑，且它的和等于 \(\operatorname{Tr}(u)\)。

只需处理 \(u\) 形如 \(\theta_{\mathrm{E}}(x^{*} \otimes x)\)（其中 \(x \in \mathrm{E}\) 且 \(x^{*} \in \mathrm{E}^{*}\)）的情形。这时族 \((\langle x, e_{i}^{*} \rangle)_{i \in \mathrm{I}}\) 有有限支撑，且我们有 \(x = \sum_{i \in \mathrm{I}} \langle x, e_{i}^{*} \rangle e_{i}\)。于是族 \((\langle x, e_{i}^{*} \rangle \langle e_{i}, x^{*} \rangle)_{i \in \mathrm{I}}\) 也有有限支撑，且我们有 \(\langle x, x^{*} \rangle = \sum_{i \in \mathrm{I}} \langle x, e_{i}^{*} \rangle \langle e_{i}, x^{*} \rangle\)。现在，我们有 \(\langle u(e_{i}), e_{i}^{*} \rangle = \langle x, e_{i}^{*} \rangle \langle e_{i}, x^{*} \rangle\)（对每个 \(i \in \mathrm{I}\)）。这就证明了该命题。

命题 2. --- 设 E、F 是投射 A-模。设 \(u \in \operatorname{Hom}_{\mathrm{A}}^{\mathrm{f}}(\mathrm{E}, \mathrm{F})\)，\(v \in \operatorname{Hom}_{\mathrm{A}}(\mathrm{F}, \mathrm{E})\)。我们有关系式

\[
\operatorname{Tr} (v \circ u) = \operatorname{Tr} (u \circ v).\tag{3}
\]

只需在 \(u\) 形如 \(\theta(x^{*} \otimes y)\)（其中 \(x^{*} \in \mathrm{E}^{*}\) 且 \(y \in \mathrm{F}\)）时证明该命题。在这种情形，我们有

\[
v \circ u = \theta_ {\mathrm{E}} (x ^ {*} \otimes v (y)) \quad \text { and } \quad u \circ v = \theta_ {\mathrm{F}} (^ {t} v (x ^ {*}) \otimes y)  ,
\]

从而

\[
\operatorname{Tr} (v \circ u) = \langle v (y), x ^ {*} \rangle = \langle y, ^ {t} v (x ^ {*}) \rangle = \operatorname{Tr} (u \circ v).
\]

推论. --- 设 E 是一个投射 A-模，u 是 \(\operatorname{End}_{\mathrm{A}}^{f}(\mathrm{E})\) 的一个元素，F 是 E 的一个包含 Im u 的投射 A-子模。用 \(u_{F}\) 表示由 u 诱导的 F 的自同态。我们有

\[
\operatorname{Tr} (u) = \operatorname{Tr} (u _ {\mathrm{F}}).\tag{4}
\]

用 i 表示 F 到 E 中的典范单射，用 \(v: E \rightarrow F\) 表示由 u 得到的同态。我们有 \(u_{F} = v \circ i\) 与 \(u = i \circ v\)。由此即得该推论。

设 E 是一个投射 A-模，\(u \in \operatorname{End}_{\mathrm{A}}^{\mathrm{f}}(\mathrm{E})\)。对每个自然数 p，A-模 \(\wedge^{p}E\) 是投射的（III, §7, No.~8, 第 519 页，推论 2），且自同态 \(\wedge^{p}u\) 属于 \(\operatorname{End}_{\mathrm{A}}^{\mathrm{f}}(\wedge^{p}\mathrm{E})\)（III, §7, No.~3, 第 511 页，命题 6），并对充分大的 p 为零。集合 \(1_{\mathrm{E}} + \operatorname{End}_{\mathrm{A}}^{\mathrm{f}}(\mathrm{E})\) 对复合是稳定的。我们定义一个从 \(1_{\mathrm{E}} + \operatorname{End}_{\mathrm{A}}^{\mathrm{f}}(\mathrm{E})\) 到 A 的映射 det，置

\[
\det (1 _ {\mathrm{E}} + u) = \sum_ {p \geqslant 0} \operatorname{Tr} \wedge^ {p} u
\]

（对 \(u\in \mathrm{End}_{\mathrm{A}}^{\mathrm{f}}(\mathrm{E})\)）。

若 E 是自由的且有限维，则由 III, §8, No.~5, 第 530 页的推论，这个定义与 III, §8, No.~1, 第 522 页中的定义一致。

命题 3. --- 设 E 是一个投射 A-模。

\begin{enumerate}
\def\labelenumi{\alph{enumi})}
\tightlist
\item
  设 \(u \in \mathrm{End}_{\mathrm{A}}^{\mathrm{f}}(\mathrm{E})\)。设 \(\mathrm{F}\) 是 \(\mathrm{E}\) 的一个包含 \(\operatorname{Im} u\) 的投射 A-子模，\(u_{\mathrm{F}}\) 是 \(\mathrm{F}\) 的由 \(u\) 诱导的自同态。我们有
\end{enumerate}

\[
\det (1 _ {\mathrm{E}} + u) = \det (1 _ {\mathrm{F}} + u _ {\mathrm{F}}).\tag{5}
\]

\begin{enumerate}
\def\labelenumi{\alph{enumi})}
\setcounter{enumi}{1}
\tightlist
\item
  设 \(u\) 与 \(v\) 是 \(\mathrm{End}_{\mathrm{A}}^{\mathrm{f}}(\mathrm{E})\) 的两个元素。我们有
\end{enumerate}

\[
\det \bigl ((1 _ {\mathrm{E}} + u) \circ (1 _ {\mathrm{E}} + v) \bigr) = \det (1 _ {\mathrm{E}} + u) \det (1 _ {\mathrm{E}} + v).\tag{6}
\]

让我们证明 a)。对每个整数 \(p \geqslant 0\)，投射 A-模 \(\wedge^p\mathrm{F}\) 可以与 \(\wedge^p\mathrm{E}\) 的一个子模等同（III, §7, No.~9, 第 520 页，推论）。\(\wedge^p u\) 的像包含在 \(\wedge^p\mathrm{F}\) 中，且 \(\wedge^p\mathrm{F}\) 的由 \(\wedge^{p}u\) 诱导的自同态等于 \(\wedge^{p}u_{F}\)。于是由命题 2 的推论，我们有 \(\operatorname{Tr}(\wedge^{p}u)=\operatorname{Tr}(\wedge^{p}u_{F})\)，从而得 a)。

让我们证明 b)。设 G 是一个使 A-模 \(L = E \oplus G\) 为自由的 A-模。用 \(u'\) 与 \(v'\) 表示 L 的自同态 \(u \oplus 0_{G}\) 与 \(v \oplus 0_{G}\)。由 a)，我们有关系式 \(\det(1_{L} + u') = \det(1_{E} + u)\)、\(\det(1_{L} + v') = \det(1_{E} + v)\)，以及

\[
\det (1 _ {\mathrm{L}} + u ^ {\prime} + v ^ {\prime} + u ^ {\prime} \circ v ^ {\prime}) = \det (1 _ {\mathrm{E}} + u + v + u \circ v).
\]

因此只需在 A-模 E 为自由时证明断言 b)。这时存在 E 的一个有限生成自由子模 F，它包含 u 的像与 v 的像。置 \(w = u + v + u \circ v\)。w 的像包含在 F 中，且我们有 \(w_{F} = u_{F} + v_{F} + u_{F} \circ v_{F}\)。于是由 (5)，我们有 \(\det(1_{E} + u) = \det(1_{F} + u_{F})\)、\(\det(1_{E} + v) = \det(1_{F} + v_{F})\)，以及

\[
\begin{array}{c} \det \bigl ((1 _ {\mathrm{E}} + u) \circ (1 _ {\mathrm{E}} + v) \bigr) \\ = \det (1 _ {\mathrm{E}} + w) = \det (1 _ {\mathrm{F}} + w _ {\mathrm{F}}) = \det \bigl ((1 _ {\mathrm{F}} + u _ {\mathrm{F}}) \circ (1 _ {\mathrm{F}} + v _ {\mathrm{F}}) \bigr). \end{array}
\]

由于 F 是一个有限生成的自由 A-模，我们有

\[
\det \left(\left(1 _ {\mathrm{F}} + u _ {\mathrm{F}}\right) \circ \left(1 _ {\mathrm{F}} + v _ {\mathrm{F}}\right)\right) = \det \left(1 _ {\mathrm{F}} + u _ {\mathrm{F}}\right) \det \left(1 _ {\mathrm{F}} + v _ {\mathrm{F}}\right) = \det \left(1 _ {\mathrm{E}} + u\right) \det \left(1 _ {\mathrm{E}} + v\right).
\]

